\section{任务与评分：Close-ended、Open-ended 与可验证性}

评估协议的第一步不是选 judge，而是判断任务输出空间和正确性来源。close-ended 与 open-ended 描述答案空间；verifiable 与 non-verifiable 描述是否存在可执行 oracle。两个维度不能混为一谈：代码生成看似开放，但可由测试验证；多项选择答案封闭，却可能因为题目歧义没有可靠真值。

\subsection{Close-ended：自动化容易，但协议仍会出错}

slides 10--14 说明 close-ended task 的候选答案和正确答案数量有限，因此可用 exact match、accuracy、precision、recall、F1 等自动评分。情感分类、事实问答和选择题是典型例子。自动化降低成本，却不保证测到了真实能力：解析器、标签定义、类别不平衡、prompt 模板和数据污染都可能改变结果。

\lecturefigure{slides-images/slide-010.png}{官方 slide 10：close/open、verifiable/non-verifiable、static/dynamic 三组维度}
\lecturefigure{slides-images/slide-011.png}{官方 slide 11：close-ended 与 open-ended 对照}
\lecturefigure{slides-images/slide-013.png}{官方 slide 13：情感与事实问答的 close-ended 示例}
\lecturefigure{slides-images/slide-014.png}{官方 slide 14：Accuracy、Precision、Recall、F1 与任务边界}

\begin{knowledgebox}{背景概念：Accuracy、Precision、Recall 与 F1}
二分类中，Precision $=TP/(TP+FP)$ 表示预测为正的样本有多少真的为正；Recall $=TP/(TP+FN)$ 表示真实正例有多少被找出；$F1=2PR/(P+R)$ 平衡 precision 与 recall。Accuracy 在类别极不平衡时可能误导，例如 99\% 样本为负时，永远预测负也有 99\% accuracy。选择指标必须对应错误成本。
\end{knowledgebox}

\begin{warningbox}{自动评分不等于客观}
Exact match 可能把语义等价答案判错，宽松字符串规则又可能接受错误答案。close-ended benchmark 还可能只覆盖易自动化的窄任务，使团队误以为 Agent 整体能力提升。评估报告应列出解析失败、歧义样本、类别分布和人工复核结果。
\end{warningbox}

\subsection{Open-ended：先区分 verifiable 与 non-verifiable}

slides 15--18 将开放任务进一步分为可验证和不可直接验证。数学证明、代码、环境状态变更虽然输出空间巨大，却可由 math verifier、unit tests 或 state oracle 检查；故事、写作和主观改写没有唯一真值，需要人类或 LLM judge 按 rubric 评价。核心问题不是“答案长不长”，而是是否存在可重复的结果判据。

\lecturefigure{slides-images/slide-015.png}{官方 slide 15：长文本、摘要、翻译与开放生成}
\lecturefigure{slides-images/slide-016.png}{官方 slide 16：开放任务分为 verifiable 与 non-verifiable}
\lecturefigure{slides-images/slide-017.png}{官方 slide 17：数学、代码和环境结果的可执行 oracle}
\lecturefigure{slides-images/slide-018.png}{官方 slide 18：故事与写作等无客观 ground truth 的任务}

\begin{importantbox}{先找 oracle，再找 judge}
若任务可通过程序、数据库约束或环境状态验证，应优先构造 deterministic oracle，并用 judge 处理解释质量等补充维度。直接用 LLM judge 替代可执行验证，会引入不必要的随机性和偏差。若没有 oracle，则必须把 rubric 拆成可观察维度，并用多评审、校准集和一致性分析控制噪声。
\end{importantbox}

\begin{knowledgebox}{读图：可验证的是 outcome，不一定是整个过程}
代码通过测试只验证已覆盖行为，未证明没有隐藏副作用；网页最终状态正确，也可能通过越权操作达到。Agent eval 常需组合 outcome verifier、trajectory constraints 和 safety policy。slide 17 给出 oracle 的入口，后面的 outcome-validity slides 会进一步讨论结果判据如何失真。
\end{knowledgebox}

\subsection{Human evaluation：金标准也有速度与一致性问题}

slides 19--20 列出人工评估维度：整体质量、流畅性、一致性、事实正确、完整性、风格、语法和冗余。人类能够理解语境和微妙偏好，因此常被视为 gold standard；但人工评估慢、昂贵，存在 inter-annotator disagreement，同一标注者随时间也可能改变标准，而且流程不易完全复现。

\lecturefigure{slides-images/slide-019.png}{官方 slide 19：人工评估的质量维度}
\lecturefigure{slides-images/slide-020.png}{官方 slide 20：人工评估的成本、一致性与复现问题}

\begin{knowledgebox}{人工评估不是“找三个人打分”}
高质量 human eval 需要任务定义、rubric、正反例、培训、盲测、随机顺序、重复样本和 disagreement adjudication。应报告标注者背景、样本数、一致性指标与置信区间。对专业法律、医疗或安全任务，普通众包偏好不能替代领域专家判断。
\end{knowledgebox}

\textbf{老师强调：}人工被称为 gold standard，是因为某些开放价值最终由人判断，不是因为人类标签无噪声。项目若把 human score 当绝对真值，会把标注偏见训练进模型。更稳健的做法是让人工负责校准、争议样本和高风险维度，让自动 oracle 承担可验证部分。

\subsection{LLM as a Judge：扩展速度，也扩展偏差}

slides 21--22 介绍用强 LLM 比较或评分候选输出。它便宜、快速、可按 rubric 扩展，且在部分任务上与人类高度相关；问题是 judge 可能不可靠、偏爱特定长度/风格、同源模型、答案位置，复杂推理还可能被表面流畅度欺骗。多次调用、chain-of-thought 或 judge ensemble 能降低方差，却不能消除系统性偏差。

\lecturefigure{slides-images/slide-021.png}{官方 slide 21：LLM judge 与人类评估的替代/补充关系}
\lecturefigure{slides-images/slide-022.png}{官方 slide 22：LLM judge 的可靠性、偏差与重复调用问题}

\begin{knowledgebox}{术语消化：Pointwise、Pairwise 与 Reference-based Judge}
Pointwise 对单个输出给分，易受绝对量表漂移；Pairwise 比较 A/B，通常更稳定但有位置偏差；reference-based judge 将候选与标准答案或 rubric 对照，适合有参考但不宜 exact match 的任务。Agent 轨迹还可做 step-wise judge，但成本高且会把 judge 的过程偏好注入评分。
\end{knowledgebox}

\begin{warningbox}{Judge 与被评模型不能形成封闭自证}
若同一模型家族生成训练数据、优化输出并担任最终 judge，系统可能学会其特定偏好。应使用可执行 oracle、不同来源 judge、人类校准和对抗样本作为外部锚点。报告中要固定 judge 版本和 prompt；在线模型更新会让历史分数失去可比性。
\end{warningbox}

\subsection{Static 与 Dynamic Eval：稳定比较和持续新鲜度的交换}

前面讨论了谁来评分，这里转向测试题本身如何维护。一个 benchmark 若永远不变，会逐渐被训练数据和开发过程吸收；若每次都变，又失去历史比较。因此本小节把稳定 anchor 与动态新题放在同一版本体系中，明确两者分别承担回归和抗污染职责。
slide 23 对比静态 benchmark 与动态 benchmark。固定测试集易复现、易做长期趋势，但会饱和、泄漏并诱导针对性优化；动态 eval 持续或周期生成新任务，降低记忆和污染风险，却增加版本控制、难度校准和跨时间比较的成本。

\lecturefigure{slides-images/slide-023.png}{官方 slide 23：static benchmark 与 dynamic benchmark}

\begin{importantbox}{双轨评估策略}
保留一个稳定 core set 做回归和历史比较，再用 rotating/dynamic set 检测新能力、污染和适应性。动态题必须保留生成器、seed、难度与 verifier 版本，必要时用 anchor items 将不同批次标尺对齐。新鲜不等于高质量，动态生成同样需要人工审计。
\end{importantbox}

\subsection{本章小结}

任务类型决定评分方法：close-ended 适合自动指标，open-ended 要先寻找 oracle，再决定 human 或 LLM judge。任何 judge 都需要校准，static 与 dynamic 则在可比性和抗污染之间权衡。下一章把这些基础维度组织成 Agent-specific taxonomy。

